\parttitle[signaling]{Cell Signaling and Signal Transduction}{Receptors, second messengers, cascades, and the logic that makes a cell respond to one molecule in a billion.}

\section{Why cells signal, the four stages, and the distance classes}

A cell is a closed chemical factory with a hydrophobic wall $3$ to $4\ \mathrm{nm}$
thick across which almost nothing polar diffuses. Everything an animal cell knows
about the outside world it learns from proteins that straddle or cross that wall, or
from the small minority of ligands lipid-soluble enough to walk through it. The whole
subject of signal transduction is the engineering answer to one problem: how does a
$180\ \mathrm{Da}$ molecule at $10^{-10}\ \mathrm{M}$ outside a cell change the
activity of an enzyme at $10^{-5}\ \mathrm{M}$ inside it, in seconds, without any of
the signal molecule ever entering?

\begin{keybox}[title={MEMORIZE: the four stages of cell signalling}]
\nr{09-signaling-four-stages}{Four stages of cell signalling}
\begin{enumerate}
\item \tm{Synthesis and release} by the signalling cell, then transport to the target.
\item \tm{Reception}: the \tm{ligand} (first messenger) binds a \tm{receptor} with
high affinity and high specificity; binding is non-covalent and reversible.
\item \tm{Transduction}: a conformational change in the receptor is converted into an
intracellular change, usually a relay of $2$ to $5$ steps involving \tm{second
messengers} and \tm{protein kinases}, with gain at each step.
\item \tm{Response}: altered enzyme activity (milliseconds to seconds), altered
cytoskeleton or transport (seconds), or altered transcription (minutes to hours).
\end{enumerate}
A fifth stage, \tm{termination}, is required for the system to be a signal at all,
and the exam treats it as part of the cycle rather than an afterthought.
\end{keybox}

Earl Sutherland's formulation is still the cleanest: reception, transduction,
response. Modern texts split synthesis off the front and termination off the back.
Both countings appear in USABO answer keys, so if a question asks "how many stages"
read the choices and pick the one that matches the list the choices describe.

\begin{thmbox}[title={PRINCIPLE: the response is a property of the target, not the signal}]
\nr{09-signaling-target-determines-response}{Target determines the response}
The same ligand produces different outcomes in different cells because the outcome is
set by the receptor subtype present, the G protein it couples to, and the downstream
effectors expressed. Epinephrine at $10^{-8}\ \mathrm{M}$ dilates skeletal muscle
arterioles ($\beta_2$, cAMP up, relaxation) and constricts skin arterioles
($\alpha_1$, \ce{IP3} and \ce{Ca^2+} up, contraction). One molecule, opposite
mechanical results, in the same animal at the same instant.
\end{thmbox}

\begin{defbox}[title={DEFINITIONS: distance classes of signalling}]
\nr{09-signaling-distance-classes}{Distance classes of signalling}
\begin{description}
\item[Intracrine] The ligand acts inside the cell that made it, never leaving it.
\item[Autocrine] The cell responds to its own secretion. Growth-factor autocrine
loops are a hallmark of tumours; activated T cells make \ce{IL-2} and carry the
\ce{IL-2} receptor.
\item[Juxtacrine (contact-dependent)] Membrane-anchored ligand on one cell binds a
receptor on a touching cell; range is the $20\ \mathrm{nm}$ of a cleft. Notch--Delta
is the archetype.
\item[Paracrine] Local diffusion, effective over $20$ to $200\ \mathrm{\mu m}$ and
$0.1$ to $10\ \mathrm{s}$, terminated by extracellular enzymes and by dilution.
\item[Synaptic] A specialised paracrine case: $20$ to $40\ \mathrm{nm}$ cleft,
transit time $\approx 0.1\ \mathrm{ms}$, cleft ligand reaching $0.1$ to
$1\ \mathrm{mM}$ transiently.
\item[Endocrine] Bloodborne, $0.1\ \mathrm{m}$ to $2\ \mathrm{m}$, latency seconds to
hours, plasma concentrations typically $10^{-12}$ to $10^{-8}\ \mathrm{M}$.
\item[Pheromonal] Between organisms, through air or water; the silkmoth responds to
bombykol at an estimated $10^{-13}\ \mathrm{g\,mL^{-1}}$ of air.
\end{description}
\end{defbox}

\begin{tabular}{@{}lllp{3.0cm}@{}}
\toprule
\textbf{Class} & \textbf{Range} & \textbf{Latency} & \textbf{Ligand concentration at target}\\
\midrule
Juxtacrine & $20\ \mathrm{nm}$ & $<1\ \mathrm{s}$ & effectively local, not diffusible\\
Synaptic & $20$--$40\ \mathrm{nm}$ & $0.1$--$2\ \mathrm{ms}$ & $10^{-4}$--$10^{-3}\ \mathrm{M}$ peak\\
Paracrine & $20$--$200\ \mathrm{\mu m}$ & $0.1$--$10\ \mathrm{s}$ & $10^{-9}$--$10^{-7}\ \mathrm{M}$\\
Endocrine & $0.1$--$2\ \mathrm{m}$ & $\mathrm{s}$--$\mathrm{h}$ & $10^{-12}$--$10^{-8}\ \mathrm{M}$\\
\bottomrule
\end{tabular}

Diffusion sets the boundary between paracrine and endocrine, and the arithmetic is
worth carrying. For a small protein with $D \approx 10^{-10}\ \mathrm{m^2\,s^{-1}}$,
the mean square displacement in three dimensions is $\langle x^2 \rangle = 6Dt$, so
$20\ \mathrm{\mu m}$ takes $t = (2\times10^{-5})^2/(6\times10^{-10}) \approx
0.7\ \mathrm{s}$, while $2\ \mathrm{mm}$ would take $6700\ \mathrm{s}$, nearly two
hours. Diffusion is fine for a cell diameter and hopeless for a body; hence a
circulation.

\begin{obsbox}[title={WORTH NOTICING: one molecule in a billion}]
\nr{09-signaling-sensitivity-numbers}{Sensitivity of signalling systems}
A hepatocyte holds roughly $10^{11}$ molecules of water and responds to epinephrine at
$10^{-10}\ \mathrm{M}$, about $6\times10^{4}$ ligand molecules in the volume
surrounding one cell, occupying perhaps $1\%$ of its $10^{4}$ receptors. A rod
photoreceptor responds detectably to a \emph{single} photon isomerising one of its
$10^{8}$ rhodopsins, a fractional occupancy of $10^{-8}$. Signalling is the biological
example of an amplifier with a gain above $10^{5}$ and a noise floor near one event.
\end{obsbox}

\begin{warnbox}[title={TRAP: reception does not have to happen at the surface}]
\nr{09-signaling-intracellular-receptors-trap}{Intracellular receptors}
The wrong belief: all receptors are membrane proteins, because signals cannot cross
membranes. The correction: steroid hormones ($\mathrm{MW}\ 250$--$500$, $\log P$
about $3$ to $5$), thyroid hormone, retinoids, vitamin D and nitric oxide all cross
the bilayer and bind receptors in the cytosol or nucleus. Their responses are slow
(transcription, $30\ \mathrm{min}$ to hours) and long-lasting (hours to days),
because the rate-limiting step is mRNA and protein synthesis, not the binding.
\end{warnbox}

\begin{tipbox}[title={TECHNIQUE: classify the signal before you classify the pathway}]
Trigger: any question naming an unfamiliar ligand. Ask three questions in order.
\begin{itemize}
\item Is it hydrophobic and small? Then predict an intracellular receptor, a
transcriptional response, and a latency of $\geq 30\ \mathrm{min}$.
\item Is it a peptide, amine or nucleotide? Then predict a surface receptor and a
second messenger, with a latency of milliseconds to minutes.
\item Does it act on a cell it touches? Then predict juxtacrine signalling with no
diffusible messenger and no amplification.
\end{itemize}
Two-thirds of multiple-choice items in this area are solved by that triage alone.
\end{tipbox}

\section{Signal classes: hormones, local regulators, neurotransmitters, gases, contact}

\begin{defbox}[title={DEFINITIONS: the chemical families of signals}]
\nr{09-signaling-ligand-chemical-classes}{Chemical classes of ligands}
\begin{description}
\item[Peptide and protein hormones] Insulin ($5808\ \mathrm{Da}$, $51$ residues, two
chains), glucagon ($29$ residues), \ce{ADH} ($9$ residues), growth hormone
($191$ residues). Stored in vesicles, released by regulated exocytosis, half-lives
$3$ to $30\ \mathrm{min}$. Water-soluble, so surface receptors only.
\item[Amine hormones and transmitters] Epinephrine, norepinephrine, dopamine,
serotonin, histamine; derived from tyrosine, tryptophan or histidine by
decarboxylation and hydroxylation. Thyroxine is the exception: an iodinated tyrosine
dimer that behaves like a steroid and rides on thyroxine-binding globulin.
\item[Steroid hormones] Cortisol, aldosterone, testosterone, estradiol, progesterone;
all from cholesterol ($27$ carbons) by side-chain cleavage. Not stored, made on
demand, $90$ to $99\%$ protein-bound in plasma, half-lives $60$ to $100\ \mathrm{min}$
(cortisol) up to $\approx 7\ \mathrm{d}$ (thyroxine).
\item[Eicosanoids] Prostaglandins, thromboxanes, leukotrienes from arachidonic acid
($20{:}4$); paracrine only, half-lives $30\ \mathrm{s}$ to a few minutes.
\item[Gases] \ce{NO} ($t_{1/2}$ $2$ to $30\ \mathrm{s}$), \ce{CO}, \ce{H2S}.
\item[Nucleotides and derivatives] \ce{ATP} and adenosine acting on P2 and P1
receptors; \ce{ATP} is a genuine extracellular transmitter at many synapses.
\end{description}
\end{defbox}

\begin{tabular}{@{}lp{2.1cm}p{2.3cm}p{2.6cm}@{}}
\toprule
\textbf{Property} & \textbf{Peptide} & \textbf{Steroid} & \textbf{Neurotransmitter}\\
\midrule
Receptor site & plasma membrane & cytosol or nucleus & membrane, post-synaptic\\
Stored? & yes, vesicles & no, made on demand & yes, vesicles\\
Plasma carrier & mostly free & $90$--$99\%$ bound & not bloodborne\\
Latency & $\mathrm{s}$--$\mathrm{min}$ & $30\ \mathrm{min}$--$\mathrm{h}$ & $0.1$--$2\ \mathrm{ms}$\\
Duration & $\mathrm{min}$ & $\mathrm{h}$--$\mathrm{d}$ & $\mathrm{ms}$\\
Amplification & large, cascade & modest, transcription & none needed, direct gating\\
\bottomrule
\end{tabular}

\begin{keybox}[title={MEMORIZE: solubility predicts everything}]
\nr{09-signaling-solubility-rule}{Solubility rule for receptor location}
Hydrophilic ligand $\rightarrow$ surface receptor $\rightarrow$ second messenger
$\rightarrow$ fast, amplified, brief response. Lipophilic ligand $\rightarrow$
intracellular receptor $\rightarrow$ direct action as a transcription factor
$\rightarrow$ slow, modest, durable response. The only common exceptions are
thyroxine (lipophilic but needs the \ce{MCT8} transporter to cross membranes at
physiological rates) and the membrane-initiated rapid actions of estradiol
(seconds, via \ce{GPER}).
\end{keybox}

\begin{expbox}[title={EXPERIMENT: Loewi 1921, the first chemical transmitter}]
\nr{09-signaling-loewi-vagusstoff}{Loewi's Vagusstoff experiment}
Two frog hearts were perfused in series. Stimulating the vagus of the first slowed it
from about $60$ to $30\ \mathrm{beats\,min^{-1}}$; the perfusate then reaching the
second, unstimulated heart slowed that one too, with a delay of a few seconds.
Physostigmine, which blocks cholinesterase, prolonged and deepened the effect;
atropine abolished it. Conclusion: nerves act by releasing a diffusible chemical,
later identified as acetylcholine, not by electrical continuity. This experiment
created the concept of a ligand and a receptor as separable objects.
\end{expbox}

\begin{obsbox}[title={WORTH NOTICING: gases break every rule}]
\nr{09-signaling-gas-messengers}{Gaseous messengers}
\ce{NO} has no receptor at the plasma membrane, no transporter, no vesicle and no
reuptake system. It is made on demand by nitric oxide synthase, diffuses roughly
$100\ \mathrm{\mu m}$ in its $2$ to $30\ \mathrm{s}$ lifetime, and binds a haem iron
inside the target cell. Its concentration is regulated purely by synthesis and decay
rate. Because it is uncharged and small ($30\ \mathrm{Da}$), it crosses membranes
faster than it diffuses laterally, and a source is therefore a sphere of influence
rather than a line of contact.
\end{obsbox}

\begin{warnbox}[title={TRAP: the same molecule is a hormone and a transmitter}]
\nr{09-signaling-same-molecule-trap}{One molecule, several signalling classes}
The wrong belief: a molecule is either a hormone or a neurotransmitter.
Norepinephrine is a neurotransmitter at sympathetic nerve terminals (cleft
concentration transiently $10^{-5}\ \mathrm{M}$, action in $\mathrm{ms}$) and a
hormone from the adrenal medulla (plasma $10^{-9}\ \mathrm{M}$, action over minutes).
Epinephrine is about $80\%$ of adrenal medullary output and norepinephrine about
$20\%$. The classification names the delivery route, not the chemical.
\end{warnbox}

Contact-dependent signalling deserves separate treatment because it has no messenger
to dilute and therefore no amplification stage. Notch on one cell binds Delta on
another; the receptor is then cleaved twice and its intracellular tail travels to the
nucleus. The stoichiometry is roughly one-to-one, and the gain is close to unity. Gap
junctions are the limiting case: connexon channels of $1.5$ to $2\ \mathrm{nm}$ pore
diameter pass anything under about $1000\ \mathrm{Da}$, including \ce{Ca^2+},
\ce{IP3} and cAMP, so coupled cells share second messengers directly. A cardiac
myocyte is electrically continuous with its neighbours through roughly
$2\times10^{5}$ connexons per intercalated disc.

\begin{tipbox}[title={TECHNIQUE: latency tells you the mechanism}]
Trigger: a data table giving response times. Map them.
\begin{itemize}
\item $0.1$--$5\ \mathrm{ms}$: ligand-gated ion channel, direct gating.
\item $0.1$--$10\ \mathrm{s}$: G protein and second messenger, or a covalent
modification of an existing protein.
\item $10$--$60\ \mathrm{min}$: transcription started, new mRNA.
\item $>2\ \mathrm{h}$: new protein accumulated, or cell division and growth.
\end{itemize}
A response abolished by cycloheximide is a transcription-or-translation response; a
response surviving it is a covalent or gating response.
\end{tipbox}

\section{Ligand-receptor binding: affinity, $K_d$, occupancy, agonists, spare receptors}

Binding is a chemical equilibrium and every quantitative question in this volume
descends from it. Write $L$ for free ligand, $R$ for free receptor, $LR$ for complex.

\begin{keybox}[title={MEMORIZE: the occupancy equation}]
\nr{09-signaling-occupancy-equation}{Receptor occupancy equation}
\[
L + R \rightleftharpoons LR, \qquad
K_d=\frac{k_{\text{off}}}{k_{\text{on}}}=\frac{[L][R]}{[LR]},\qquad
Y=\frac{[LR]}{[R]_{\text{total}}}=\frac{[L]}{K_d+[L]}
\]
$Y$ is fractional occupancy (dimensionless), $[L]$ free ligand concentration in
$\mathrm{M}$, $K_d$ the dissociation constant in $\mathrm{M}$, $k_{\text{on}}$ in
$\mathrm{M^{-1}\,s^{-1}}$ (diffusion limit about $10^{8}$ to
$10^{9}\ \mathrm{M^{-1}\,s^{-1}}$) and $k_{\text{off}}$ in $\mathrm{s^{-1}}$. The
equation assumes one site per receptor, no cooperativity, and $[L]$ not depleted by
binding. At $[L]=K_d$, $Y=0.5$ exactly. Going from $Y=0.1$ to $Y=0.9$ costs an
$81$-fold rise in $[L]$ for any non-cooperative receptor.
\end{keybox}

\begin{defbox}[title={DEFINITIONS: the pharmacology vocabulary}]
\nr{09-signaling-pharmacology-terms}{Agonist, antagonist, efficacy, potency}
\begin{description}
\item[Affinity] How tightly the ligand binds; measured by $K_d$ (or $K_i$ for a
competitor). Low $K_d$ means high affinity. Typical hormone receptors:
$K_d = 10^{-11}$ to $10^{-8}\ \mathrm{M}$; typical neurotransmitter receptors at a
synapse: $10^{-6}$ to $10^{-4}\ \mathrm{M}$, because the cleft concentration is high
and speed matters more than affinity.
\item[Efficacy] The size of the maximal response the ligand can produce. A
\tm{full agonist} gives the tissue maximum, a \tm{partial agonist} gives $10$ to
$80\%$ of it however high the dose, an \tm{antagonist} gives zero, an
\tm{inverse agonist} drives the response below the unstimulated baseline.
\item[Potency] The concentration needed, reported as $\mathrm{EC_{50}}$. Set by
affinity, receptor number, and efficiency of coupling. Potency and efficacy are
independent axes.
\item[Competitive antagonist] Binds the orthosteric site reversibly, shifts the
dose-response curve right with no fall in maximum; surmountable.
\item[Non-competitive or allosteric antagonist] Binds elsewhere or irreversibly,
lowers the maximum; insurmountable.
\end{description}
\end{defbox}

\begin{calcbox}[title={WORKED: occupancy, potency and a shifted curve}]
\nr{09-signaling-occupancy-worked}{Worked occupancy and antagonist shift}
An insulin receptor has $K_d = 1.0\times10^{-10}\ \mathrm{M}$. Fasting plasma insulin
is $5\ \mathrm{\mu U\,mL^{-1}}$, which is $3.0\times10^{-11}\ \mathrm{M}$; after a
meal it reaches $60\ \mathrm{\mu U\,mL^{-1}}$, or
$3.6\times10^{-10}\ \mathrm{M}$.

\ttitle{Fasting occupancy.}
\[
Y=\frac{3.0\times10^{-11}}{1.0\times10^{-10}+3.0\times10^{-11}}
=\frac{3.0}{13.0}=0.231
\]
\ttitle{Fed occupancy.}
\[
Y=\frac{3.6\times10^{-10}}{1.0\times10^{-10}+3.6\times10^{-10}}
=\frac{3.6}{4.6}=0.783
\]
So a $12$-fold rise in hormone changes occupancy from $23\%$ to $78\%$, a factor of
only $3.4$. Hyperbolic binding compresses a wide input range into a narrow output
range: that compression is exactly why cascades must supply gain downstream.

\ttitle{Now add a competitive antagonist.} A drug with $K_i=2.0\times10^{-9}\ \mathrm{M}$
is present at $1.8\times10^{-8}\ \mathrm{M}$. The dose-ratio is
\[
r=1+\frac{[I]}{K_i}=1+\frac{1.8\times10^{-8}}{2.0\times10^{-9}}=1+9=10
\]
so the apparent $K_d$ becomes $1.0\times10^{-9}\ \mathrm{M}$ and fed occupancy falls
to $3.6/(10+3.6)=0.265$.

\ttitle{Check.} At $[L]=K_d$ the formula must give $0.5$: with
$[L]=1.0\times10^{-10}$, $Y=1.0/(1.0+1.0)=0.500$. \tm{$\surd$} And the antagonist
must not change the maximum: as $[L]\rightarrow\infty$, $Y\rightarrow 1$ with or
without the drug. \tm{$\surd$}
\end{calcbox}

\begin{thmbox}[title={PRINCIPLE: spare receptors uncouple $\mathrm{EC_{50}}$ from $K_d$}]
\nr{09-signaling-spare-receptors}{Spare receptors}
If maximal response requires occupancy of only a fraction of receptors, the tissue has
\tm{spare receptors} (receptor reserve) and $\mathrm{EC_{50}} < K_d$. In smooth
muscle, full contraction can be reached at $5$ to $10\%$ occupancy, so
$\mathrm{EC_{50}}$ may be $10$ to $20$ times lower than $K_d$. Two consequences the
exam tests: a partial irreversible block of receptors first shifts the curve left-of-
$K_d$ region rightwards without lowering the maximum, and only once the reserve is
consumed does the maximum fall; and a tissue with spare receptors responds faster,
because the time to reach threshold occupancy is shorter.
\end{thmbox}

\begin{tabular}{@{}lp{2.4cm}p{2.3cm}p{2.5cm}@{}}
\toprule
\textbf{Ligand type} & \textbf{Max response} & \textbf{Curve shift} & \textbf{Example}\\
\midrule
Full agonist & $100\%$ & reference & epinephrine at $\beta_1$\\
Partial agonist & $10$--$80\%$ & own curve & buprenorphine at $\mu$\\
Competitive antagonist & unchanged & right, parallel & propranolol at $\beta$\\
Non-competitive & lowered & little shift & ketamine at \ce{NMDA}\\
Inverse agonist & below basal & --- & many \ce{H1} blockers\\
\bottomrule
\end{tabular}

\begin{tipbox}[title={TECHNIQUE: read a binding plot in three seconds}]
\nr{09-signaling-binding-plots}{Binding plot geometry}
On a linear plot, saturation binding is a rectangular hyperbola: read $B_{\max}$ off
the plateau and $K_d$ at half-plateau. On a semi-log plot the same data is a sigmoid
with the steepest point at $\mathrm{EC_{50}}$, and the useful width between $10\%$
and $90\%$ is $1.91$ log units for a non-cooperative site. On a Scatchard plot,
$B/[L]$ against $B$, the slope is $-1/K_d$ and the $x$-intercept is $B_{\max}$. A
curved, concave-up Scatchard plot means two affinity classes or negative
cooperativity; the insulin receptor genuinely shows this.
\end{tipbox}

\begin{warnbox}[title={TRAP: $K_d$ is not $\mathrm{EC_{50}}$ and neither is $K_m$}]
\nr{09-signaling-kd-ec50-trap}{Distinguishing $K_d$, $\mathrm{EC_{50}}$ and $K_m$}
The wrong belief: $K_d$, $\mathrm{EC_{50}}$ and $K_m$ are interchangeable
half-maximal constants. $K_d$ is a binding equilibrium constant, and lower means
tighter. $\mathrm{EC_{50}}$ is a functional constant that also contains receptor
density, amplification and desensitisation, and it can be $100$-fold below $K_d$ in a
tissue with spare receptors. $K_m$ is a steady-state kinetic constant for catalysis
and only equals a true $K_d$ when $k_{\text{cat}} \ll k_{\text{off}}$. A question
giving you both $K_d$ and $\mathrm{EC_{50}}$ is almost always testing receptor
reserve.
\end{warnbox}

\section{G protein coupled receptors: structure, the cycle, the $\mathrm{G}\alpha$ families}

\begin{keybox}[title={MEMORIZE: GPCR architecture and scale}]
\nr{09-signaling-gpcr-architecture}{GPCR architecture}
A single polypeptide of $300$ to $1100$ residues threads the membrane seven times as
$\alpha$ helices of $20$ to $27$ residues each, giving an extracellular N-terminus, an
intracellular C-terminus, three extracellular loops and three intracellular loops.
The ligand pocket sits in the upper third of the helical bundle for small amines, or
on the N-terminal domain and loops for peptides and glycoprotein hormones. Loop
$\mathrm{IC}3$ and the C-terminal tail carry the G protein contact and the
phosphorylation sites. The human genome encodes about $800$ GPCRs, of which roughly
$400$ are odorant receptors; GPCRs are the direct target of $30$ to $35\%$ of
approved drugs.
\end{keybox}

Agonist binding rotates helix $6$ outward by about $10$ angstroms at its
cytoplasmic end and shifts the conserved \ce{DRY} motif, opening a cleft that the
G protein's $\alpha$ subunit docks into. The receptor is a catalyst: it is a
guanine nucleotide exchange factor, and one occupied receptor can activate many
G proteins in sequence. That is the first amplification step in the cascade.

\begin{defbox}[title={DEFINITIONS: the heterotrimeric G protein}]
\nr{09-signaling-g-protein-subunits}{G protein subunits}
\begin{description}
\item[$\mathrm{G}\alpha$] $39$ to $52\ \mathrm{kDa}$, binds GDP or GTP, has intrinsic
GTPase activity with $k_{\text{cat}}$ of roughly $2$ to
$5\ \mathrm{min^{-1}}$ alone. Lipid-anchored by palmitoylation or myristoylation.
\item[$\mathrm{G}\beta\gamma$] An obligate dimer, $35$ and $8\ \mathrm{kDa}$;
$\gamma$ is prenylated. It is a signalling molecule in its own right, gating
\ce{GIRK} \ce{K+} channels, inhibiting \ce{Ca_V} channels, and activating
\ce{PLC}$\beta$ and \ce{PI3K}$\gamma$.
\item[RGS proteins] Regulators of G protein signalling; GTPase-activating proteins
that raise $\mathrm{G}\alpha$ hydrolysis up to $2000$-fold, shortening the active
state from tens of seconds to well under a second.
\end{description}
Humans have $16$ $\mathrm{G}\alpha$ genes, $5$ $\beta$ and $12$ $\gamma$, allowing on
the order of $10^{3}$ combinations.
\end{defbox}

\begin{thmbox}[title={PRINCIPLE: the G protein cycle is a timer, not a switch}]
\nr{09-signaling-g-protein-cycle}{The G protein cycle}
\begin{enumerate}
\item Resting: $\mathrm{G}\alpha$-GDP is bound to $\mathrm{G}\beta\gamma$; the
trimer is inactive and membrane-tethered.
\item Agonist-bound receptor acts as a GEF: GDP leaves, and cytosolic GTP
($200$ to $500\ \mathrm{\mu M}$, about $10$-fold above GDP) binds.
\item GTP binding reorders three switch regions; $\mathrm{G}\alpha$-GTP loses
affinity for $\beta\gamma$ and both halves engage effectors.
\item Intrinsic GTPase, accelerated by RGS, hydrolyses GTP to GDP $+$ \ce{Pi}. The
lifetime of the active state, typically $1$ to $30\ \mathrm{s}$ without RGS and
$50\ \mathrm{ms}$ to $1\ \mathrm{s}$ with it, is what sets the duration of the
response.
\item $\mathrm{G}\alpha$-GDP reassociates with $\beta\gamma$.
\end{enumerate}
Nothing is phosphorylated in this cycle, and no external energy is consumed beyond
the one GTP. The receptor does not hold the G protein on; the clock in
$\mathrm{G}\alpha$ does.
\end{thmbox}

\begin{tabular}{@{}lp{2.6cm}p{2.6cm}p{2.3cm}@{}}
\toprule
\textbf{Family} & \textbf{Effector} & \textbf{Second messenger} & \textbf{Toxin target}\\
\midrule
$\mathrm{G}\alpha_s$ & adenylyl cyclase up & cAMP up & cholera toxin\\
$\mathrm{G}\alpha_{i/o}$ & adenylyl cyclase down & cAMP down & pertussis toxin\\
$\mathrm{G}\alpha_q$ & \ce{PLC}$\beta$ up & \ce{IP3}, DAG, \ce{Ca^2+} & none common\\
$\mathrm{G}\alpha_{12/13}$ & RhoGEF up & Rho-GTP, actin & none common\\
$\mathrm{G}\alpha_t$ (transducin) & \ce{cGMP} PDE6 up & cGMP down & cholera-like\\
\bottomrule
\end{tabular}

\begin{obsbox}[title={WORTH NOTICING: transducin is $\mathrm{G}\alpha_i$ wearing a hat}]
\nr{09-signaling-transducin}{Transducin as a $\mathrm{G}\alpha_i$ relative}
Transducin belongs to the $i/o$ family and, like them, lowers a cyclic nucleotide,
except that it does so by activating a phosphodiesterase rather than inhibiting a
cyclase. The logic is identical: agonist (here, a photon absorbed by retinal)
$\rightarrow$ GPCR (rhodopsin) $\rightarrow$ G protein $\rightarrow$ change in cyclic
nucleotide $\rightarrow$ change in a channel. Vision is a chemical signalling cascade
that happens to be triggered by light, and rhodopsin's ligand, $11$-\emph{cis}
retinal, is covalently bound until the photon isomerises it.
\end{obsbox}

\begin{expbox}[title={EXPERIMENT: Gilman's cyc$^-$ complementation, 1977--1980}]
\nr{09-signaling-gilman-cyc-minus}{Gilman's cyc$^-$ experiment}
\ce{S49} mouse lymphoma mutants called cyc$^-$ had normal $\beta$-adrenergic binding
sites and normal adenylyl cyclase activity when assayed with detergent and
\ce{Mn^2+}, yet made no cAMP in response to isoproterenol. Mixing cyc$^-$ membranes
with a soluble fraction from wild-type cells restored hormone-stimulated cyclase.
Gilman purified the restoring activity and found a GTP-binding protein of $45$ and
$35\ \mathrm{kDa}$ subunits. Rodbell had already shown that GTP, not \ce{ATP}, was
required for hormonal activation and that the non-hydrolysable analogue
\ce{GTP}$\gamma$\ce{S} gave a persistent response. Together these established that a
separate transducer sits between receptor and enzyme, which is the central idea of
the field. Shared Nobel Prize, $1994$.
\end{expbox}

\begin{warnbox}[title={TRAP: three errors about G proteins}]
\nr{09-signaling-g-protein-traps}{Common G protein errors}
The wrong beliefs and their corrections.
\begin{itemize}
\item "GTP phosphorylates the G protein." No: GTP \emph{binds}. Nothing gains a
phosphate. The receptor is a nucleotide exchanger, not a kinase.
\item "$\mathrm{G}\beta\gamma$ is just an inhibitory chaperone." No: free
$\beta\gamma$ opens \ce{GIRK} channels in the heart, which is how acetylcholine at
\ce{M2} receptors slows the sinoatrial node within $100\ \mathrm{ms}$, far faster than
any cAMP change.
\item "Hydrolysis of GTP activates the G protein." The reverse: hydrolysis
\emph{terminates} it. A GTPase-deficient $\mathrm{G}\alpha_s$ mutant, as in
McCune-Albright syndrome and some pituitary adenomas, is constitutively active.
\end{itemize}
\end{warnbox}

\section{The cyclic AMP pathway, adenylyl cyclase, PKA, and amplification arithmetic}

\begin{keybox}[title={MEMORIZE: the canonical cAMP cascade}]
\nr{09-signaling-camp-cascade}{The cAMP cascade}
Epinephrine $\rightarrow$ $\beta_2$ receptor $\rightarrow$ $\mathrm{G}\alpha_s$-GTP
$\rightarrow$ adenylyl cyclase $\rightarrow$ cAMP $\rightarrow$ protein kinase A
$\rightarrow$ phosphorylase kinase $\rightarrow$ glycogen phosphorylase
$\rightarrow$ glucose $1$-phosphate. In parallel, PKA phosphorylates glycogen
synthase and inactivates it, so one signal turns synthesis off and breakdown on.
The reaction made by the cyclase is
\[
\ce{ATP -> cAMP + PPi}
\]
a cyclisation between the $3'$-\ce{OH} and the $\alpha$-phosphate, with
pyrophosphate released and then hydrolysed, which makes the step effectively
irreversible.
\end{keybox}

Adenylyl cyclase is a $120\ \mathrm{kDa}$ integral protein with twelve
transmembrane helices in two bundles and two cytoplasmic catalytic domains, C1 and
C2, that form one active site at their interface. Nine membrane isoforms plus a
soluble bicarbonate-sensitive one exist in mammals; \ce{AC1} and \ce{AC8} are
stimulated by \ce{Ca^2+}-calmodulin, \ce{AC5} and \ce{AC6} are inhibited by
\ce{Ca^2+}, and all membrane isoforms need \ce{Mg^2+} and are activated by the
diterpene forskolin. That isoform dependence is how a cell decides whether the
calcium and cAMP pathways reinforce or oppose one another.

\begin{defbox}[title={DEFINITIONS: protein kinase A}]
\nr{09-signaling-pka-structure}{Protein kinase A structure and activation}
Inactive PKA is a heterotetramer \ce{R2C2}: two regulatory and two catalytic
subunits. Each R subunit has two tandem cAMP sites, \ce{CNB-A} and \ce{CNB-B}, and
binding is positively cooperative with a Hill coefficient near $1.6$ to $1.8$ and a
half-maximal activation at $100$ to $500\ \mathrm{nM}$ cAMP. Four cAMP molecules
release two free C subunits of $40\ \mathrm{kDa}$. The C subunit phosphorylates
serine and threonine in the consensus \ce{R-R-X-S/T-\Phi}, where $\Phi$ is a bulky
hydrophobic residue. \tm{PKI} is a heat-stable inhibitor peptide; \tm{AKAPs} are
scaffold proteins that tether PKA within $50\ \mathrm{nm}$ of chosen substrates.
Resting cytosolic cAMP is about $0.05$ to $0.1\ \mathrm{\mu M}$ and stimulated cAMP
reaches $1$ to $10\ \mathrm{\mu M}$, so the operating range straddles the PKA
$\mathrm{K_{1/2}}$ almost exactly.
\end{defbox}

\begin{expbox}[title={EXPERIMENT: Sutherland and Rall, 1956--1958}]
\nr{09-signaling-sutherland-camp}{Sutherland's discovery of cAMP}
Epinephrine activated glycogen phosphorylase in intact liver slices, but adding
epinephrine to a soluble liver extract did nothing. Sutherland and Rall found that
epinephrine \emph{plus a particulate membrane fraction plus \ce{ATP} and \ce{Mg^2+}}
generated a small, heat-stable, acid-stable factor that, transferred to the soluble
fraction on its own, activated phosphorylase with no hormone present. The factor was
identified as adenosine $3',5'$-cyclic monophosphate. Three conclusions defined the
field: the hormone never enters the cell; an intracellular \tm{second messenger}
carries the message; and the messenger is generated at the membrane and acts in the
cytosol. Nobel Prize, $1971$. The same logic later identified cGMP, \ce{IP3} and
\ce{Ca^2+}: separate the compartments, show the transferable factor.
\end{expbox}

\begin{calcbox}[title={WORKED: the amplification arithmetic}]
\nr{09-signaling-amplification-arithmetic}{Amplification arithmetic of the cAMP cascade}
A hepatocyte carries about $2\times10^{4}$ $\beta$ receptors. Suppose $1\%$ are
occupied, so $200$ active receptors, and take these per-step gains: each active
receptor activates $10$ \ce{Gs} during its lifetime; each active cyclase makes
$1000$ cAMP; each PKA catalytic subunit phosphorylates $10$ phosphorylase kinase
molecules; each phosphorylase kinase activates $50$ phosphorylase; each phosphorylase
releases $10^{4}$ glucose $1$-phosphate.

\ttitle{Step by step.}
\begin{itemize}
\item Active \ce{Gs} and hence cyclase: $200 \times 10 = 2\times10^{3}$.
\item cAMP molecules: $2\times10^{3} \times 10^{3} = 2\times10^{6}$.
\item PKA activated, taking $4$ cAMP per \ce{R2C2} and $2$ C subunits released:
$2\times10^{6}/4 = 5\times10^{5}$ tetramers, giving $10^{6}$ free C subunits.
\item Phosphorylase kinase: $10^{6}\times 10 = 10^{7}$.
\item Phosphorylase: $10^{7}\times 50 = 5\times10^{8}$.
\item Glucose $1$-phosphate: $5\times10^{8}\times10^{4}=5\times10^{12}$ molecules.
\end{itemize}
\ttitle{Convert to mass.}
$5\times10^{12}/6.02\times10^{23}=8.3\times10^{-12}\ \mathrm{mol}$; at
$180\ \mathrm{g\,mol^{-1}}$ for glucose that is $1.5\times10^{-9}\ \mathrm{g}$, about
$1.5\ \mathrm{ng}$ of glucose from one cell.

\ttitle{Overall gain.} $5\times10^{12}$ products from $200$ bound ligands is
$2.5\times10^{10}$-fold, and from the $200$ hormone molecules that is a gain of about
$10^{10}$.

\ttitle{Check two ways.} Multiply the gains: $10\times10^{3}\times0.5\times10\times50
\times10^{4}=2.5\times10^{10}$, matching the step-by-step total. \tm{$\surd$} And
sanity-check the physiology: a hepatocyte holds about $200\ \mathrm{pg}$ of glycogen,
so $1.5\ \mathrm{ng}$ would exceed the store; real cells stop far short because
phosphodiesterases, phosphatases and receptor desensitisation truncate the cascade
within seconds. The arithmetic gives the \emph{capacity}, not the yield. \tm{$\surd$}
\end{calcbox}

\begin{tabular}{@{}lp{2.5cm}ll@{}}
\toprule
\textbf{Step} & \textbf{Species} & \textbf{Gain} & \textbf{Cumulative}\\
\midrule
$1$ & receptor--\ce{Gs} & $\times 10$ & $10^{1}$\\
$2$ & cyclase--cAMP & $\times 10^{3}$ & $10^{4}$\\
$3$ & cAMP--PKA C & $\times 0.5$ & $5\times10^{3}$\\
$4$ & PKA--phos.\ kinase & $\times 10$ & $5\times10^{4}$\\
$5$ & kinase--phosphorylase & $\times 50$ & $2.5\times10^{6}$\\
$6$ & phosphorylase--G1P & $\times 10^{4}$ & $2.5\times10^{10}$\\
\bottomrule
\end{tabular}

Notice that step $3$ has a gain below one. Cascades are not uniformly amplifying;
they trade gain for control at the steps where a switch is wanted, and the
cooperativity of cAMP binding to PKA converts a graded cAMP rise into a sharper
kinase response.

\begin{obsbox}[title={WORTH NOTICING: cAMP is ancient and bacterial}]
\nr{09-signaling-camp-cap}{cAMP in bacteria}
\sn{Escherichia coli} uses cAMP with no receptor and no G protein at all: glucose
starvation relieves inhibition of adenylyl cyclase, cAMP rises about $50$-fold to
roughly $1\ \mathrm{mM}$ locally, and cAMP-\ce{CAP} binds promoter DNA to activate
catabolic operons. Eukaryotes kept the messenger and bolted a receptor and a
transducer onto the front of it. The same molecule, opposite logic: in the bacterium
cAMP means "no glucose", in the hepatocyte it means "release glucose".
\end{obsbox}

\begin{warnbox}[title={TRAP: cAMP does not always mean PKA, or stimulation}]
\nr{09-signaling-camp-traps}{cAMP misconceptions}
Three corrections. First, cAMP has targets besides PKA: \ce{HCN} pacemaker channels
and \ce{CNG} channels bind it directly, and \ce{Epac} is a cAMP-activated RapGEF, so
a cAMP response surviving PKA inhibition by \ce{H89} is not evidence against cAMP.
Second, epinephrine does not always raise cAMP: at $\alpha_2$ receptors it couples to
$\mathrm{G}\alpha_i$ and lowers it. Third, cAMP itself is not the activator of PKA's
catalytic subunit in the sense of binding it; cAMP binds the \emph{regulatory}
subunit and relieves inhibition. Disinhibition and activation look identical on a
graph and are different mechanisms.
\end{warnbox}

\begin{tipbox}[title={TECHNIQUE: locate the lesion in a cascade}]
\nr{09-signaling-lesion-localisation}{Locating a cascade lesion}
Trigger: a table of treatments and responses. Add each downstream agent in turn and
ask what rescues the response.
\begin{itemize}
\item No response to hormone, normal response to forskolin: lesion is at the
receptor or the G protein, not the cyclase.
\item No response to hormone or forskolin, normal response to
dibutyryl-cAMP (membrane-permeant): lesion is the cyclase.
\item No response to dibutyryl-cAMP, normal response to phorbol ester: lesion is PKA
or below, and the PKC arm is intact.
\end{itemize}
The principle generalises: the first agent that rescues the response acts
\emph{downstream} of the lesion.
\end{tipbox}

\section{The phosphoinositide pathway: PLC, \ce{IP3}, DAG, PKC and calcium release}

\begin{keybox}[title={MEMORIZE: one lipid, two messengers}]
\nr{09-signaling-pip2-cleavage}{Cleavage of \ce{PIP2}}
Phospholipase C cleaves phosphatidylinositol $4,5$-bisphosphate:
\[
\ce{PIP2 -> IP3 + DAG}
\]
\ce{PIP2} is only $0.2$ to $1\%$ of membrane phospholipid, roughly
$5\times10^{3}$ to $10^{4}$ molecules per $\mathrm{\mu m^2}$ of inner leaflet.
Hydrolysis is at the glycerophosphate bond, so the \emph{water-soluble} inositol
$1,4,5$-trisphosphate diffuses into the cytosol ($D\approx 2.8\times
10^{-10}\ \mathrm{m^2\,s^{-1}}$, range tens of micrometres in its $\approx1\ \mathrm{s}$
lifetime) while the \emph{lipid} diacylglycerol stays in the membrane and recruits
protein kinase C to it. The bifurcation is the whole point: one enzyme, one
substrate, two messengers with different compartments and different targets.
\end{keybox}

\begin{defbox}[title={DEFINITIONS: the phosphoinositides and their enzymes}]
\nr{09-signaling-phosphoinositide-species}{Phosphoinositide species}
\begin{description}
\item[\ce{PI}] Phosphatidylinositol, the unphosphorylated parent, about $10\%$ of
membrane phospholipid.
\item[\ce{PI(4,5)P2}] Made by \ce{PI4K} then \ce{PIP5K}; substrate of PLC and of
\ce{PI3K}. Also a direct cofactor for many channels and transporters.
\item[\ce{PI(3,4,5)P3}] Made by \ce{PI3K} acting on \ce{PIP2}; a docking site for
\ce{PH} domains, notably \ce{Akt} and \ce{PDK1}. Removed by the phosphatase
\ce{PTEN}, a tumour suppressor.
\item[\ce{PLC}$\beta$] Activated by $\mathrm{G}\alpha_q$ and by
$\mathrm{G}\beta\gamma$; the GPCR arm.
\item[\ce{PLC}$\gamma$] Has \ce{SH2} domains, activated by receptor tyrosine kinase
phosphotyrosines; the RTK arm.
\item[\ce{PLC}$\delta,\ \epsilon,\ \eta,\ \zeta$] \ce{Ca^2+}-sensitive, Ras-activated,
neuronal and sperm-specific forms respectively. All isoforms require \ce{Ca^2+} for
catalysis, which makes the pathway self-amplifying.
\end{description}
\end{defbox}

\begin{thmbox}[title={PRINCIPLE: the two branches have separate jobs}]
\nr{09-signaling-ip3-dag-branches}{The \ce{IP3} and DAG branches}
The \ce{IP3} branch is the \emph{mobilising} branch: \ce{IP3} binds the \ce{IP3}
receptor, a tetrameric \ce{Ca^2+} channel of the endoplasmic reticulum membrane with
a conductance near $10$ to $20\ \mathrm{pS}$ and a half-maximal activating \ce{IP3}
concentration of $0.1$ to $1\ \mathrm{\mu M}$, and \ce{Ca^2+} floods from a store
held at $100$ to $800\ \mathrm{\mu M}$ into a cytosol at $100\ \mathrm{nM}$. The DAG
branch is the \emph{membrane-recruiting} branch: DAG plus \ce{Ca^2+} brings
conventional PKC to the inner leaflet, where its pseudosubstrate is displaced and it
phosphorylates serine and threonine residues on membrane-proximal substrates. The
branches converge because conventional PKC needs the \ce{Ca^2+} that the other branch
released: the pathway implements a coincidence requirement inside itself.
\end{thmbox}

\begin{tabular}{@{}lp{2.4cm}p{2.3cm}p{2.5cm}@{}}
\toprule
\textbf{Feature} & \textbf{\ce{IP3}} & \textbf{DAG} & \textbf{\ce{Ca^2+}}\\
\midrule
Compartment & cytosol & membrane & cytosol, local\\
Target & \ce{IP3}R on ER & PKC, \ce{TRPC} & calmodulin, PKC, many\\
Resting level & $<0.1\ \mathrm{\mu M}$ & low & $50$--$100\ \mathrm{nM}$\\
Peak level & $1$--$5\ \mathrm{\mu M}$ & $5$--$10\times$ basal & $0.5$--$10\ \mathrm{\mu M}$\\
Removal & $5$-phosphatase, $3$-kinase & DAG kinase to \ce{PA}, lipase & \ce{SERCA}, \ce{PMCA}, exchanger\\
Lifetime & $\approx 1\ \mathrm{s}$ & $\mathrm{s}$--$\mathrm{min}$ & $\mathrm{ms}$--$\mathrm{s}$ locally\\
\bottomrule
\end{tabular}

\begin{calcbox}[title={WORKED: how many calcium ions does one \ce{IP3} release?}]
\nr{09-signaling-calcium-flux-worked}{Calcium released per \ce{IP3} receptor opening}
An \ce{IP3} receptor channel has a single-channel \ce{Ca^2+} current of about
$0.5\ \mathrm{pA}$ and stays open for roughly $10\ \mathrm{ms}$ per event.

\ttitle{Charge moved.}
\[
Q = I t = 0.5\times10^{-12}\ \mathrm{A}\times 10\times10^{-3}\ \mathrm{s}
= 5\times10^{-15}\ \mathrm{C}
\]
\ttitle{Ions.} \ce{Ca^2+} carries $2$ elementary charges, so
\[
n=\frac{5\times10^{-15}}{2\times1.602\times10^{-19}}
=\frac{5\times10^{-15}}{3.204\times10^{-19}}
\approx 1.56\times10^{4}\ \text{ions}
\]
\ttitle{Concentration change in a microdomain.} In a hemisphere of radius
$0.5\ \mathrm{\mu m}$, volume
$\tfrac{2}{3}\pi r^{3}=\tfrac{2}{3}\pi(0.5\times10^{-6})^{3}
=2.6\times10^{-19}\ \mathrm{m^3}=2.6\times10^{-16}\ \mathrm{L}$, so
\[
\Delta c=\frac{1.56\times10^{4}/6.02\times10^{23}}{2.6\times10^{-16}}
=\frac{2.6\times10^{-20}}{2.6\times10^{-16}}=1.0\times10^{-4}\ \mathrm{M}
\]
that is $100\ \mathrm{\mu M}$ in the microdomain, a thousandfold above resting
cytosolic \ce{Ca^2+}.

\ttitle{Check.} Cells buffer $98$ to $99\%$ of entering \ce{Ca^2+}, so the free
concentration actually seen is $1$ to $2\%$ of $100\ \mathrm{\mu M}$, i.e.\ $1$ to
$2\ \mathrm{\mu M}$, which is precisely the measured amplitude of a \ce{Ca^2+} puff.
\tm{$\surd$} Whole-cell averaging over a $1\ \mathrm{pL}$ cell would give
$\Delta c=2.6\times10^{-20}/10^{-12}=26\ \mathrm{nM}$, invisible: the signal exists
only because it is local. \tm{$\surd$}
\end{calcbox}

\begin{obsbox}[title={WORTH NOTICING: the PKC family is three families}]
\nr{09-signaling-pkc-isoforms}{PKC isoform classes}
Conventional PKC ($\alpha,\ \beta\mathrm{I},\ \beta\mathrm{II},\ \gamma$) needs DAG
and \ce{Ca^2+} and phosphatidylserine. Novel PKC ($\delta,\ \epsilon,\ \eta,\
\theta$) needs DAG but not \ce{Ca^2+}, because its \ce{C2} domain lacks the
calcium-coordinating aspartates. Atypical PKC ($\zeta,\ \iota/\lambda$) needs
neither, and is regulated instead by \ce{PIP3} and protein partners. Nine to eleven
genes, one catalytic chemistry, three completely different input requirements: a
single kinase family implementing three different logic gates.
\end{obsbox}

\begin{warnbox}[title={TRAP: five errors in the phosphoinositide pathway}]
\nr{09-signaling-pip-traps}{Phosphoinositide pathway traps}
\begin{itemize}
\item "\ce{IP3} opens channels in the plasma membrane." No: the \ce{IP3} receptor is
an \emph{ER} channel. Plasma membrane \ce{Ca^2+} entry in this pathway is
store-operated, through \ce{Orai1} gated by \ce{STIM1}, and is a consequence of store
depletion, not of \ce{IP3}.
\item "DAG diffuses to the nucleus." No: DAG has two fatty acyl chains and stays in
the bilayer. That is why PKC must come to the membrane.
\item "\ce{IP3} and \ce{PIP3} are the same thing." They are not even the same class:
\ce{IP3} is a soluble sugar phosphate from PLC; \ce{PIP3} is a membrane lipid from
\ce{PI3K}. One mobilises calcium, the other recruits \ce{Akt}.
\item "Phorbol esters are PKC inhibitors." They are potent, poorly metabolised PKC
\emph{activators} that substitute for DAG, which is why they are tumour promoters.
\item "$\mathrm{G}\alpha_q$ pathways do not use cAMP so the two are independent."
Cross-talk is heavy: \ce{Ca^2+}-calmodulin activates \ce{AC1} and \ce{AC8} and
inhibits \ce{AC5} and \ce{AC6}, and PKC phosphorylates and modulates several cyclases.
\end{itemize}
\end{warnbox}

\begin{tipbox}[title={TECHNIQUE: tell a $\mathrm{G}_q$ response from a $\mathrm{G}_s$ response}]
\nr{09-signaling-gq-vs-gs-discrimination}{Discriminating $\mathrm{G}_q$ from $\mathrm{G}_s$}
Trigger: a question that gives a hormone and a measured output.
\begin{itemize}
\item Response abolished in \ce{Ca^2+}-free medium containing \ce{EGTA} and
unaffected by forskolin: $\mathrm{G}_q$.
\item Response mimicked by forskolin or by dibutyryl-cAMP: $\mathrm{G}_s$.
\item Response mimicked by a phorbol ester: the DAG-PKC arm of $\mathrm{G}_q$.
\item Response abolished by pertussis toxin: $\mathrm{G}_{i/o}$, and note this is the
one family the toxin reports on.
\item Smooth muscle contracting rather than relaxing: think $\mathrm{G}_q$ and
\ce{Ca^2+}-calmodulin-myosin light chain kinase, and check
\pref{physiology} for the mechanics.
\end{itemize}
\end{tipbox}
